\documentclass[12pt]{article}
\usepackage[utf8]{inputenc}
\usepackage[T1]{fontenc}
\usepackage[margin=1in]{geometry}
\usepackage{setspace}
\usepackage{xcolor}
\usepackage[colorlinks=true, urlcolor=blue!60!black, linkcolor=blue!60!black]{hyperref}
\setlength{\parindent}{2em}
\setlength{\parskip}{1.25em}
\renewcommand{\baselinestretch}{1.0}
\title{\textit{Zelman v. Simmons-Harris}, 536 U.S. 639 (2002)}
\author{Case Brief}
\date{}
\begin{document}
\maketitle

% ============================================================================
% Constitutional Law Briefing Book --- Chapter: Religious Liberty
% Case Brief: Zelman v. Simmons-Harris, 536 U.S. 639 (2002)
% ============================================================================

\section{Case Metadata}

\begin{description}
  \item[Case Name:] \emph{Zelman, Superintendent of Public Instruction of Ohio, et al.\ v.\ Simmons-Harris et al.} (consolidated with \emph{Hanna Perkins School v.\ Simmons-Harris} and \emph{Taylor v.\ Simmons-Harris})
  \item[Citation:] \href{https://supreme.justia.com/cases/federal/us/536/639/}{536 U.S. 639 (2002)}
  \item[Court:] Supreme Court of the United States (on certiorari to the U.S. Court of Appeals for the Sixth Circuit, 234 F.3d 945 (6th Cir. 2000))
  \item[Argued:] February 20, 2002 \quad \textbf{Decided:} June 27, 2002
  \item[Vote:] 5--4
  \item[Majority Opinion:] Chief Justice William H. Rehnquist (joined by Justices O'Connor, Scalia, Kennedy, and Thomas)
  \item[Concurring Opinions:] Justice Sandra Day O'Connor; Justice Clarence Thomas
  \item[Dissenting Opinions:] Justice John Paul Stevens; Justice David H. Souter (joined by Justices Stevens, Ginsburg, and Breyer); Justice Stephen G. Breyer (joined by Justices Stevens and Souter)
  \item[Provisions at Issue:] U.S. Const. amend.~I (Establishment Clause), incorporated via amend.~XIV; Ohio Pilot Project Scholarship Program, Ohio Rev. Code Ann. \S\S~3313.974--3313.979
  \item[Chapter Placement:] Religious Liberty
\end{description}

\section{Facts and Procedural History}

In 1995, a federal district court declared a ``crisis of magnitude'' and placed the Cleveland City School District---serving over 75,000 predominantly low-income and minority students---under state control after the district failed all 18 state minimal performance standards and saw over two-thirds of its high school students drop or fail out before graduation (\href{https://supreme.justia.com/cases/federal/us/536/639/}{536 U.S. at 644}).

In response, Ohio enacted the Pilot Project Scholarship Program for any district under federal-court-ordered state supervision (affecting only Cleveland) (\href{https://supreme.justia.com/cases/federal/us/536/639/}{536 U.S. at 644--645}). The program offered two forms of aid: (1) tuition scholarships for K--8 students to attend participating private schools (religious or nonreligious) within Cleveland or public schools in adjacent districts, and (2) tutorial grants (90\% up to \$360 for low-income families; 75\% for others) for students remaining in Cleveland public schools (\href{https://supreme.justia.com/cases/federal/us/536/639/}{536 U.S. at 645--647}). Low-income families (below 200\% of the poverty line) received priority for 90\% of tuition up to \$2,250, with parental copayments capped at \$250; other families received 75\% up to \$1,875. Checks were payable to parents, who endorsed them to the chosen school. Participating private schools had to meet state educational standards and agree not to discriminate on the basis of race, religion, or ethnicity, or teach hatred (\href{https://supreme.justia.com/cases/federal/us/536/639/}{536 U.S. at 645--646}).

During the 1999--2000 school year, 56 private schools participated, 46 (82\%) of which were religiously affiliated; no adjacent suburban public schools participated; and 96\% of the more than 3,700 scholarship students enrolled in religious schools (\href{https://supreme.justia.com/cases/federal/us/536/639/}{536 U.S. at 647}). Meanwhile, Cleveland also offered 10 nonreligious publicly funded ``community'' (charter) schools (enrolling over 1,900 students at \$4,518 per pupil) and 23 public magnet schools (enrolling over 13,000 students at \$7,746 per pupil) (\href{https://supreme.justia.com/cases/federal/us/536/639/}{536 U.S. at 647--648}).

After the Ohio Supreme Court rejected federal Establishment Clause claims but invalidated the statute on state procedural grounds---which the legislature promptly cured---Ohio taxpayers sued in federal court in July 1999 (\href{https://supreme.justia.com/cases/federal/us/536/639/}{536 U.S. at 648}). The District Court granted summary judgment for the taxpayers (72 F. Supp. 2d 834), and a divided Sixth Circuit panel affirmed (234 F.3d 945), holding that the program had the ``primary effect'' of advancing religion. The Supreme Court granted certiorari and reversed.

\section{Issue(s) Presented}

Does Ohio's Pilot Project Scholarship Program---which provides religion-neutral tuition vouchers to Cleveland families that may be redeemed at religious or secular private schools, and under which 96\% of participating scholarship recipients enroll in religious schools---violate the Establishment Clause of the First Amendment?

\section{Holding}

\textbf{No (5--4).} A government aid program that is neutral with respect to religion and provides assistance directly to a broad class of citizens who, in turn, direct the aid to religious schools wholly as a result of their own genuine and independent private choice does not violate the Establishment Clause (\href{https://supreme.justia.com/cases/federal/us/536/639/}{536 U.S. at 652--653, 662--663}).

\section{Reasoning}

Chief Justice Rehnquist applied the ``purpose'' and ``effect'' framework of \emph{Agostini v. Felton}, 521 U.S. 203 (1997) (\href{https://supreme.justia.com/cases/federal/us/536/639/}{536 U.S. at 648--649}; \href{https://www.law.cornell.edu/constitution-conan/amendment-1/zelman-and-indirect-aid}{Cornell LII}). Because the program's secular purpose---assisting poor children in a failing school district---was undisputed, the case turned solely on whether the program had the forbidden ``effect'' of advancing religion.

\begin{itemize}
  \item \textbf{Direct Aid vs.\ True Private Choice:} Relying on \emph{Mueller v. Allen}, 463 U.S. 388 (1983), \emph{Witters v. Washington Dep't of Services for the Blind}, 474 U.S. 481 (1986), and \emph{Zobrest v. Catalina Foothills School Dist.}, 509 U.S. 1 (1993), the Court drew a sharp line between direct aid to religious schools and ``programs of true private choice'' (\href{https://supreme.justia.com/cases/federal/us/536/639/}{536 U.S. at 649--653}). When public funds reach religious schools only through the independent decisions of numerous private individuals, ``the circuit between government and religion [is] broken,'' and any advancement or perceived endorsement of religion is attributable to individual parents, not the State (\href{https://supreme.justia.com/cases/federal/us/536/639/}{536 U.S. at 652--655}).
  \item \textbf{Neutrality and Absence of Financial Skewing:} Ohio's program conferred benefits on a broad class defined purely by residence and financial need, permitted all district private schools and adjacent public schools to participate, and actually created financial \emph{disincentives} against choosing private religious schools: private schools received half the per-pupil state aid of community schools and one-third that of magnet schools, while families had to pay a copayment at private schools but nothing at public, community, or magnet schools (\href{https://supreme.justia.com/cases/federal/us/536/639/}{536 U.S. at 653--654}).
  \item \textbf{Evaluating the Full Menu of Educational Options:} Rejecting respondents' reliance on the 82\% religious-school and 96\% religious-enrollment figures, the Court held that genuine secular choice must be assessed across \emph{all} options Ohio offered Cleveland children---including traditional public schools, public schools with tutoring aid, nonreligious private schools, community schools, and magnet schools (\href{https://supreme.justia.com/cases/federal/us/536/639/}{536 U.S. at 655--660}). Including community, magnet, and tutored public students reduced religious enrollment to under 20\% (\href{https://supreme.justia.com/cases/federal/us/536/639/}{536 U.S. at 659}). Tying constitutionality to annual enrollment tallies would yield absurd geographic disparities across cities and States (\href{https://supreme.justia.com/cases/federal/us/536/639/}{536 U.S. at 657--658}).
  \item \textbf{Distinguishing \emph{Nyquist}:} The Court distinguished \emph{Committee for Public Education \& Religious Liberty v. Nyquist}, 413 U.S. 756 (1973), which invalidated New York tuition grants restricted exclusively to private-school parents, holding that \emph{Nyquist} does not govern neutral programs offering aid broadly without regard to religion (\href{https://supreme.justia.com/cases/federal/us/536/639/}{536 U.S. at 661--662}).
\end{itemize}

\section{Concurring and Dissenting Opinions}

\begin{itemize}
  \item \textbf{Justice O'Connor (Concurring):} Emphasized that the decision did not mark a sharp break from precedent or longstanding indirect subsidies (tax exemptions, Pell Grants, the G.I. Bill, Medicare/Medicaid) and defended counting Cleveland's community and magnet schools as genuine nonreligious options (\href{https://supreme.justia.com/cases/federal/us/536/639/}{536 U.S. at 663--676}).
  \item \textbf{Justice Thomas (Concurring):} Questioned whether the Establishment Clause---originally a federalism restraint protecting States from a national church---should bind States identically through the Fourteenth Amendment, and invoked Frederick Douglass to argue that school choice provides vital educational emancipation for urban minority children (\href{https://supreme.justia.com/cases/federal/us/536/639/}{536 U.S. at 676--684}).
  \item \textbf{Justice Stevens (Dissenting):} Argued that Cleveland's educational crisis, the availability of public-school alternatives, and parental voluntariness were constitutionally irrelevant to whether the State may fund religious indoctrination, warning that removing bricks from the wall of separation invites religious strife (\href{https://supreme.justia.com/cases/federal/us/536/639/}{536 U.S. at 684--686}).
  \item \textbf{Justice Souter (Dissenting, joined by Stevens, Ginsburg, and Breyer):} Charged the majority with abandoning \emph{Everson v. Board of Education}'s no-aid principle in favor of ``verbal formalism'' (\href{https://supreme.justia.com/cases/federal/us/536/639/}{536 U.S. at 686--689}). Comparing vouchers to public-school spending gutted neutrality and choice as limiting principles; because the \$2,500 tuition cap favored parish-subsidized Catholic schools and few secular private seats existed, 96.6\% of voucher recipients faced a ``Hobson's choice,'' while taxpayer funding threatened liberty of conscience, state regulation of religious curricula, and sectarian discord (\href{https://supreme.justia.com/cases/federal/us/536/639/}{536 U.S. at 696--717}).
  \item \textbf{Justice Breyer (Dissenting, joined by Stevens and Souter):} Focused on the Establishment Clause's anti-divisiveness purpose in a religiously diverse Nation, warning that state funding criteria (such as Ohio's ban on teaching ``hatred'') would entangle state officials in adjudicating religious doctrine and spark sectarian conflict (\href{https://supreme.justia.com/cases/federal/us/536/639/}{536 U.S. at 717--725}).
\end{itemize}

\section{Key Quotations}

\begin{enumerate}
  \item \textbf{Chief Justice Rehnquist on True Private Choice} (\href{https://supreme.justia.com/cases/federal/us/536/639/}{536 U.S. at 652}):
  \begin{quote}
    ``[W]here a government aid program is neutral with respect to religion, and provides assistance directly to a broad class of citizens who, in turn, direct government aid to religious schools wholly as a result of their own genuine and independent private choice, the program is not readily subject to challenge under the Establishment Clause.''
  \end{quote}

  \item \textbf{Chief Justice Rehnquist on Statistical Tallies} (\href{https://supreme.justia.com/cases/federal/us/536/639/}{536 U.S. at 658}):
  \begin{quote}
    ``The constitutionality of a neutral educational aid program simply does not turn on whether and why, in a particular area, at a particular time, most private schools are run by religious organizations, or most recipients choose to use the aid at a religious school.''
  \end{quote}

  \item \textbf{Justice Thomas (Concurring) on Educational Choice and the Fourteenth Amendment} (\href{https://supreme.justia.com/cases/federal/us/536/639/}{536 U.S. at 680}):
  \begin{quote}
    ``There would be a tragic irony in converting the Fourteenth Amendment's guarantee of individual liberty into a prohibition on the exercise of educational choice.''
  \end{quote}

  \item \textbf{Justice Souter (Dissenting) on ``Hobson's Choice''} (\href{https://supreme.justia.com/cases/federal/us/536/639/}{536 U.S. at 707}):
  \begin{quote}
    ``The criterion is one of genuinely free choice on the part of the private individuals who choose, and a Hobson's choice is not a choice, whatever the reason for being Hobsonian.''
  \end{quote}
\end{enumerate}

\section{Doctrinal Significance and Modern Legacy (Through 2026)}

\emph{Zelman} transformed Religion Clauses doctrine from strict separationism to a neutrality-and-private-choice paradigm---and ultimately paved the way for a dramatic shift from \emph{permitting} religious school aid under the Establishment Clause to \emph{requiring} equal inclusion under the Free Exercise Clause:

\begin{itemize}
  \item \textbf{The Narrow ``Play in the Joints'' (\emph{Locke v. Davey}):} Two years after \emph{Zelman}, the Court held in \href{https://supreme.justia.com/cases/federal/us/540/712/}{\emph{Locke v. Davey}, 540 U.S. 712 (2004)}, that while \emph{Zelman} permitted a State to fund religious instruction via neutral scholarships, Washington did not violate the Free Exercise Clause by excluding devotional theology degrees from a state college scholarship program.
  \item \textbf{From Permitted to Required (\emph{Trinity Lutheran}, \emph{Espinoza}, and \emph{Carson}):} Over the next two decades, the Court sharply narrowed \emph{Locke} to ministerial training and converted \emph{Zelman}'s Establishment Clause shield into a Free Exercise sword. In \href{https://supreme.justia.com/cases/federal/us/582/15-577/}{\emph{Trinity Lutheran Church of Columbia, Inc. v. Comer}, 582 U.S. 449 (2017)}, the Court held that excluding a church preschool from a neutral playground-resurfacing grant based on religious status violated the Free Exercise Clause. In \href{https://supreme.justia.com/cases/federal/us/591/18-1195/}{\emph{Espinoza v. Montana Dep't of Revenue}, 591 U.S. 464 (2020)}, the Court invalidated Montana's application of a state ``Blaine Amendment'' no-aid clause to bar religious schools from a tax-credit scholarship program, citing \emph{Zelman} to confirm that indirect aid posed no Establishment Clause bar (\href{https://www.law.cornell.edu/constitution-conan/amendment-1/zelman-and-indirect-aid}{Cornell LII}). Finally, in \href{https://supreme.justia.com/cases/federal/us/596/20-1088/}{\emph{Carson v. Makin}, 596 U.S. \_\_\_ (2022)}, the Court struck down Maine's exclusion of ``sectarian'' schools from its rural tuition-assistance program, abolishing any status-versus-use distinction: once a State subsidizes private education, it cannot exclude religious schools because of their religious status \emph{or} religious instruction. That same Term, \href{https://supreme.justia.com/cases/federal/us/597/21-418/}{\emph{Kennedy v. Bremerton School District}, 597 U.S. \_\_\_ (2022)}, formally abandoned the \emph{Lemon} test in favor of historical practices and understandings.
  \item \textbf{The Public-Charter Frontier (2025--2026):} The outer boundary of \emph{Zelman} and \emph{Carson} reached the Court in \href{https://www.scotusblog.com/cases/oklahoma-statewide-charter-school-board-v-drummond/}{\emph{Oklahoma Statewide Charter School Board v. Drummond}} (consolidated with \emph{St. Isidore of Seville Catholic Virtual School v. Drummond}), testing whether a State must allow a privately run \emph{religious public charter school} under \emph{Carson} or whether charter schools are state actors barred from religious instruction. On May 22, 2025, an equally divided Court (4--4, with Justice Barrett recused) affirmed the Oklahoma Supreme Court's judgment blocking the religious charter school without a precedential opinion, leaving the boundary between private voucher schools and public charter schools unresolved through 2026.
\end{itemize}

\section{Pedagogical Critiques and Exam Hypotheticals}

\subsection{Pedagogical Critiques}
\begin{itemize}
  \item \textbf{The Denominator Problem:} \emph{Zelman}'s outcome hinged on defining the baseline of ``genuine secular choice.'' Including public magnet and community schools dropped religious participation from 96\% to under 20\%, whereas Justice Souter's voucher-only denominator exposed severe market concentration. Students should recognize how framing the baseline often dictates constitutional outcomes.
  \item \textbf{The Vanishing ``Play in the Joints'':} Together, \emph{Zelman} (neutral vouchers are permitted) and \emph{Carson} (excluding religious schools from private-school aid is forbidden) leave States with a binary choice: fund only public schools, or fund private secular and religious schools equally.
\end{itemize}

\subsection{Exam Hypothetical and ``Opposite Side'' Contrast}
\textbf{Hypothetical:} A State establishes an ``Education Savings Account'' (ESA) program granting \$7,000 per child to any family in low-performing districts to spend on private-school tuition, tutoring, or homeschooling. In one rural county, the only accredited private schools are three evangelical Christian academies, so 100\% of county ESA tuition funds flow to religious schools. Taxpayers sue under the Establishment Clause.

\textbf{Analysis:} Under \href{https://supreme.justia.com/cases/federal/us/536/639/}{\emph{Zelman}}, the ESA program survives. It is facially neutral, aids a broad class defined without reference to religion, and routes funds solely via independent parental choice. As Chief Justice Rehnquist held, constitutionality does not turn on local market demographics or annual tallies (\href{https://supreme.justia.com/cases/federal/us/536/639/}{536 U.S. at 657--658}), and families retain secular options (public schools, tutoring, secular homeschooling).

\textbf{``Opposite Side'' Contrast:} Now flip the challenge: suppose the State, fearing Establishment Clause liability in that rural county, amends the ESA statute to exclude any school that ``integrates religious doctrine into its curriculum,'' while still funding secular private schools and tutoring. Under \href{https://supreme.justia.com/cases/federal/us/591/18-1195/}{\emph{Espinoza}} and \href{https://supreme.justia.com/cases/federal/us/596/20-1088/}{\emph{Carson v. Makin}}, that exclusion violates the Free Exercise Clause. Because \emph{Zelman} establishes that neutral private-choice aid does not offend the Establishment Clause, the State has no compelling anti-establishment interest to justify discriminating against religious schools based on either religious status or religious use. By contrast, if the State instead funds \emph{only} traditional public schools and declines to offer any private-school aid program at all, the Free Exercise Clause is not violated (\href{https://supreme.justia.com/cases/federal/us/596/20-1088/}{\emph{Carson}, 596 U.S. \_\_\_ (2022)}).

\end{document}
